\input{../tex-inputs/latex-leseansicht-vorspann}

\section[Arthur Schnitzler an Rainer Maria Rilke, 13. 4. 1899]{L04341 Arthur Schnitzler an Rainer Maria Rilke, 13. 4. 1899}
\nopagebreak\mylabel{L04341v}
\rehead{ }\normalsize\beginnumbering\briefempfaengerindex{Rilke, Rainer Maria@\textsc{Rilke, Rainer Maria}!zzzSchnitzler, Arthur@\emph{von Arthur Schnitzler}!1899-04-131@{13. 4. 1899}|(be}
\toendnotes[C]{\smallbreak\pagebreak[2]}
\correspDesc{Versand  durch Arthur Schnitzler am 13. 4. 1899 in Wien
\newline{}Erhalt  durch Rainer Maria Rilke im Zeitraum [14. 4. 1899 – 18. 4. 1899?] in Prag}\toendnotes[C]{\smallbreak}
\Standort{DLA, A:Schnitzler, HS.2022.81.}
\physDesc{Briefkarte, 602 Zeichen
\newline{}Handschrift: schwarze Tinte, deutsche Kurrent}
\buchAbdrucke{\weitereDrucke{1) \emph{Rainer Maria Rilke und Arthur Schnitzler. Ihr Briefwechsel.}. Mit Anmerkungen herausgegeben von Heinrich Schnitzler In: \emph{Wort und Wahrheit}, Jg. 13, Nr. 4, April 1958, S. 282–298, hier S. 284.} \weitereDrucke{2) Arthur Schnitzler: \emph{Briefe 1875–1912}. Herausgegeben von Therese Nickl und Heinrich Schnitzler. Frankfurt am Main: \emph{S. Fischer} 1981, S. 370.} }\toendnotes[C]{\smallbreak}
\pstart
           \noindent{}{\pb}Lieber Herr Rilke, von Tag zu Tag habe ich eine Sti{\geminationm}ung abgewartet, um Ihnen \label{K_L04341-1v}\edtext{für Ihre Bücher\pwindex{Rilke, Rainer Maria 4.\,12.\,1875 Prag – 29.\,12.\,1926 Montreux@\textsc{Rilke, Rainer Maria} (4.\,12.\,1875 Prag – 29.\,12.\,1926 Montreux)!Zwei Prager Geschichten@\strich\emph{Zwei Prager Geschichten}|pwv}\pwindex{Rilke, Rainer Maria 4.\,12.\,1875 Prag – 29.\,12.\,1926 Montreux@\textsc{Rilke, Rainer Maria} (4.\,12.\,1875 Prag – 29.\,12.\,1926 Montreux)!Am Leben hin. Novellen und Skizzen@\strich\emph{Am Leben hin. Novellen und Skizzen}|pwv} herzlich zu danken}{\lemma{\textnormal{\emph{für … danken}}}\Cendnote{\textnormal{Die Übergabe
                        von \emph{Zwei Prager Geschichten}\pwindex{Rilke, Rainer Maria 4.\,12.\,1875 Prag – 29.\,12.\,1926 Montreux@\textsc{Rilke, Rainer Maria} (4.\,12.\,1875 Prag – 29.\,12.\,1926 Montreux)!Zwei Prager Geschichten@\strich\emph{Zwei Prager Geschichten}|pwk} und \emph{Am Leben hin. Novellen und Skizzen}\pwindex{Rilke, Rainer Maria 4.\,12.\,1875 Prag – 29.\,12.\,1926 Montreux@\textsc{Rilke, Rainer Maria} (4.\,12.\,1875 Prag – 29.\,12.\,1926 Montreux)!Am Leben hin. Novellen und Skizzen@\strich\emph{Am Leben hin. Novellen und Skizzen}|pwk} 
                        dürfte am 16. 3. 1899 beim persönlichen Zusammentreffen
                  stattgefunden haben. Das Vorwort von \emph{Zwei Prager Geschichten}\pwindex{Rilke, Rainer Maria 4.\,12.\,1875 Prag – 29.\,12.\,1926 Montreux@\textsc{Rilke, Rainer Maria} (4.\,12.\,1875 Prag – 29.\,12.\,1926 Montreux)!Zwei Prager Geschichten@\strich\emph{Zwei Prager Geschichten}|pwk} ist mit »Februar 1899« datiert, 
                  der Druck bei 
                  \emph{Adolf Bonz {\kaufmannsund} Comp.}\orgindex{Adolf Bonz und Comp.@Adolf Bonz {\kaufmannsund}  Comp.|pwk} also gerade fertiggestellt.}}}\label{K_L04341-1} – um Ihnen mehr darüber zu{ }ſagen, besonders über das{ }ſchöne Skizzen u.
                  Novellen Buch\pwindex{Rilke, Rainer Maria 4.\,12.\,1875 Prag – 29.\,12.\,1926 Montreux@\textsc{Rilke, Rainer Maria} (4.\,12.\,1875 Prag – 29.\,12.\,1926 Montreux)!Am Leben hin. Novellen und Skizzen@\strich\emph{Am Leben hin. Novellen und Skizzen}|pwv}. Aber ich ka{\geminationn} noch immer nicht{ }ſo{ }ſchreiben als ich wollte; de{\geminationn} ein unendlich \label{K_L04341-2v}\edtext{ſchwerer Verluſt\pwindex{Reinhard, Marie 13.\,3.\,1871 Wien – 18.\,3.\,1899 ebd.@\textsc{Reinhard, Marie} (13.\,3.\,1871 Wien – 18.\,3.\,1899 ebd.)|pwv}}{\lemma{\textnormal{\emph{schwerer Verlust}}}\Cendnote{\textnormal{Seine Partnerin Marie Reinhard\pwindex{Reinhard, Marie 13.\,3.\,1871 Wien – 18.\,3.\,1899 ebd.@\textsc{Reinhard, Marie} (13.\,3.\,1871 Wien – 18.\,3.\,1899 ebd.)|pwk} war am 18. 3. 1899 an einer
                  Blutvergiftung gestorben.}}}\label{K_L04341-2}, den ich (\label{K_L04341-3v}\edtext{am
               Tag nachdem ich Sie zuletzt geſehn}{\lemma{\textnormal{\emph{am … gesehn}}}\Cendnote{\textnormal{Das datiert das letzte Treffen auf den 17. 3. 1899, an 
                  dem Rilke\pwindex{Rilke, Rainer Maria 4.\,12.\,1875 Prag – 29.\,12.\,1926 Montreux@\textsc{Rilke, Rainer Maria} (4.\,12.\,1875 Prag – 29.\,12.\,1926 Montreux)|pwk} nicht im \emph{Tagebuch}\pwindex{Schnitzler, Arthur 15. 5. 1862 Wien – 21. 10. 1931 ebd.@\textsc{Schnitzler, Arthur} (15. 5. 1862 Wien – 21. 10. 1931 ebd.)!Tagebuch@\strich\emph{Tagebuch}|pwk} Erwähnung fand. Mit einiger Wahrscheinlichkeit dürfte das
                  Treffen im Zuge der Generalprobe von \emph{Der Abenteurer und die Sängerin}\pwindex{Hofmannsthal, Hugo von 1.\,2.\,1874 Wien – 15.\,7.\,1929 Rodaun@\textsc{Hofmannsthal, Hugo von} (1.\,2.\,1874 Wien – 15.\,7.\,1929 Rodaun)!Abenteurer und die Sängerin oder Die Geschenke des Lebens@\strich\emph{Der Abenteurer und die Sängerin oder Die Geschenke des Lebens}|pwk} und \emph{Die Hochzeit der Sobeide}\pwindex{Hofmannsthal, Hugo von 1.\,2.\,1874 Wien – 15.\,7.\,1929 Rodaun@\textsc{Hofmannsthal, Hugo von} (1.\,2.\,1874 Wien – 15.\,7.\,1929 Rodaun)!Hochzeit der Sobeide@\strich\emph{Die Hochzeit der Sobeide}|pwk}\eventindex{Burgtheater@\textbf{Burgtheater}!Wiener Generalprobe von Der Abenteurer und die Sängerin und Die Hochzeit der Sobeide, 17.3.1899@Wiener Generalprobe von Der Abenteurer und die Sängerin und Die Hochzeit der Sobeide, 17.3.1899|pwk} im Burgtheater\oindex{Wien@\textbf{Wien}!I., Innere Stadt@\textbf{I., Innere Stadt}!Burgtheater@\textbf{Burgtheater}|pwk} stattgefunden haben.}}}\label{K_L04341-3}) erlitt, macht mich zu allem
               unfähig, was geordnete Gedanken vorausſetzt. Bitte entſchuldigen Sie mich alſo und
                  {\pb}glauben Sie an meine herzliche Sympathie. Vielleicht
               \label{K_L04341-4v}\edtext{begegnen wir einander nächſte Woche}{\lemma{\textnormal{\emph{begegnen … Woche}}}\Cendnote{\textnormal{Das nächste Treffen fand erst
                  dreizehn Jahre später statt, am 15. 5. 1912.}}}\label{K_L04341-4} in Berlin\oindex{Berlin@\textbf{Berlin}|pw}.\pend
           
\pstart
           \label{K_L04341-5v}\edtext{Grüßen Sie Frau \textsc{Lou Andrea Salomé\pwindex{Andreas-Salomé, Lou 12.\,2.\,1861 Sankt Petersburg – 5.\,2.\,1937 Göttingen@\textsc{Andreas-Salomé, Lou} (12.\,2.\,1861 Sankt Petersburg – 5.\,2.\,1937 Göttingen)|pw}}}{\lemma{\textnormal{\emph{Grüßen … Salomé}}}\Cendnote{\textnormal{Rilke\pwindex{Rilke, Rainer Maria 4.\,12.\,1875 Prag – 29.\,12.\,1926 Montreux@\textsc{Rilke, Rainer Maria} (4.\,12.\,1875 Prag – 29.\,12.\,1926 Montreux)|pwk} wohnte zwischen 1898 und 
                     1900 bei Lou Andreas-Salomé\pwindex{Andreas-Salomé, Lou 12.\,2.\,1861 Sankt Petersburg – 5.\,2.\,1937 Göttingen@\textsc{Andreas-Salomé, Lou} (12.\,2.\,1861 Sankt Petersburg – 5.\,2.\,1937 Göttingen)|pwk} in der
                     Villa Waldfrieden\oindex{Villa Waldfrieden@\textbf{Villa Waldfrieden}|pwk} in Schmargendorf\oindex{Schmargendorf@\textbf{Schmargendorf}|pwk}.
                  }}}\label{K_L04341-5}.{\\[\baselineskip]}Ihr \spacefill\mbox{Arthur Schnitzler}\pend
           \leftskip=0em{}
\pstart
           Wien\oindex{Wien@\textbf{Wien}|pw}{ }13. 4. 99.\pend
           \selectlanguage{ngerman}\endnumbering\briefempfaengerindex{Rilke, Rainer Maria@\textsc{Rilke, Rainer Maria}!zzzSchnitzler, Arthur@\emph{von Arthur Schnitzler}!1899-04-131@{13. 4. 1899}|)be}\mylabel{L04341h}  \newcommand{\dateiname}{L04341}\newcommand{\titel}{Arthur Schnitzler an Rainer Maria Rilke, 13. 4. 1899}\newcommand{\editorInnen}{Herausgegeben von Jahnke, SelmaMüller, Martin Anton}\input{../tex-inputs/latex-leseansicht-abspann}
